%==============================================================================
%  01-executive-summary.tex — Executive Summary
%==============================================================================
\strategyPart{01 \textbar\ Executive Summary}{The Strategic Imperative}

The core philosophy that governs every argument in this document is stated
upfront and without compromise:

\begin{center}
  \yBanner[color=inNavy]{We do not implement AI for the organization; we solve organizational problems using AI.}
\end{center}

\section{Why this document exists}

Artificial intelligence has moved from laboratory curiosity to strategic
necessity. The macroeconomic evidence is unambiguous: the McKinsey Global
Institute projects that AI could add up to \textbf{\$13~trillion} to global
economic output by 2030, while McKinsey's 2023 research estimates that
generative AI alone could contribute \textbf{\$2.6--\$4.4~trillion} in
annual value across 63 use cases [McKinsey Global Institute 2018; McKinsey
2023]. Yet the same body of research shows that most organizations capture
almost none of this value. In the MIT Sloan--BCG survey, seven of ten
companies report minimal or no impact from AI to date [MIT SMR--BCG 2019].

The gap between the promise and the outcome is not a technology gap. It is
an \textbf{organizational} gap: how problems are defined, how data and talent
are governed, how work is redesigned, and how the organization itself is
reinvented. This document provides the framework for closing that gap.

\section{Key findings at a glance}

\begin{center}
  \stat{88\%}{Organizations use AI in\\at least one function [McKinsey 2025]}\hfill
  \stat{7\%}{Report AI fully scaled\\across the organization [McKinsey 2025]}\hfill
  \stat{30\%}{Gen-AI projects abandoned\\after proof of concept [Gartner 2023]}
\end{center}

\begin{itemize}
  \item \textbf{Adoption has never been higher — scaling has never been rarer.}
        McKinsey's 2025 survey reports that \textbf{88\% of organizations} use
        AI in at least one business function, yet only \textbf{7\%} report that
        AI is fully scaled across the organization [McKinsey State of AI 2025].
  \item \textbf{Adoption is rising faster than value.} 65\% of organizations now
        regularly use generative AI in at least one function, up from roughly a
        third a year earlier [McKinsey State of AI 2024].
  \item \textbf{Failure is the default.} Gartner warned that through 2022,
        \textbf{85\% of AI projects} would deliver erroneous outcomes due to
        bias in data, algorithms, or the teams managing them [Gartner 2018];
        a further \textbf{30\%} of generative-AI projects will be abandoned
        after proof of concept through 2025 [Gartner 2023].
  \item \textbf{Value follows problems, not tools.} Organizations that anchor
        AI to a defined business problem and a measurable KPI are the ones that
        scale; the rest stall at proof of concept [Gartner 2023; McKinsey 2023].
  \item \textbf{Measured gains exist and are concentrated.} Controlled research
        shows \textbf{55.8\% faster} coding tasks with assistants [GitHub
        2022], \textbf{+14\%} customer-support productivity and \textbf{+34\%}
        for novice agents [MIT/NBER 2023], and 25--40\% faster knowledge work
        [BCG 2023].
  \item \textbf{Readiness is organizational, not technical.} The binding
        constraints are human capital, data capital, and organizational capital
        — the three ``stocks'' an AI-enabled organization must accumulate
        [MIT CISR; HBR 2018].
\end{itemize}

\section{How to read this document}

The survey is organized into five acts.

\begin{itemize}
  \item \textbf{Act 1 · The Problem} (Parts 01--02) — the strategic
        imperative and the research scope.
  \item \textbf{Act 2 · The Landscape} (Parts 03--11) — definitions, theory,
        the paradigm, transformation models, the value ladder, technologies,
        how AI works, AI by function, and data \& technology foundations.
  \item \textbf{Act 3 · The Organization} (Parts 12--17) — operating models,
        readiness, work redesign, people, capabilities and governance.
  \item \textbf{Act 4 · The Journey} (Parts 18--21) — why AI fails, case
        studies, measurement and the roadmap.
  \item \textbf{Act 5 · The Path Forward} (Parts 22--24) — comparative
        analysis, predictions and the executive conclusion.
  \item \textbf{Back matter} (Parts 25--28) — methodology, references,
        glossary and the self-assessment toolkit.
\end{itemize}

Every part ends with a \textbf{key takeaway} box; Part 28 turns the survey
into a working diagnostic for your organization.

\takeaway{Adoption is high but scaling is rare (88\% use AI, 7\% have scaled). The gap between the two is organizational, not technical.}
